
\documentclass[a4paper,10pt]{article}

\usepackage[utf8]{inputenc}
\usepackage[T1]{fontenc}
\usepackage[turkish]{babel}
\usepackage{geometry}
\usepackage{booktabs}
\usepackage{array}
\usepackage{tabularx}
\usepackage{longtable}
\usepackage{enumitem}
\usepackage{listings}
\usepackage{xcolor}
\usepackage[hidelinks]{hyperref}
\usepackage{microtype}
\usepackage{amsmath}
\usepackage{titlesec}

% Türkçe babel paketinin "=" karakterini listings seçeneklerinde etkilemesini engeller
\AtBeginDocument{\shorthandoff{=}}

\geometry{
    a4paper,
    top=1.8cm,
    bottom=1.8cm,
    left=1.8cm,
    right=1.8cm
}

\setlength{\parindent}{0pt}
\setlength{\parskip}{5pt}
\renewcommand{\arraystretch}{1.18}

\titleformat{\section}{\large\bfseries}{\thesection.}{0.5em}{}
\titleformat{\subsection}{\normalsize\bfseries}{\thesubsection.}{0.5em}{}

\lstdefinestyle{cstyle}{
    language={[ANSI]C},
    basicstyle=\ttfamily\small,
    keywordstyle=\bfseries,
    commentstyle=\itshape,
    numbers=left,
    numberstyle=\tiny,
    numbersep=7pt,
    frame=single,
    breaklines=true,
    showstringspaces=false,
    tabsize=4
}

\newcommand{\hex}[1]{\texttt{0x#1}}
\newcommand{\code}[1]{\texttt{#1}}

\title{\textbf{DEOS CAN-FD Interface Control Document (ICD)}\\
\large Sürüm 0.3}
\author{DEOS -- Dokuz Eylül Otonom Sistemler}
\date{}

\begin{document}
\maketitle

\begin{center}
\begin{tabular}{ll}
\textbf{Doküman tipi:} & Interface Control Document (ICD) \\
\textbf{Protokol sürümü:} & v1.1 (\hex{11}) \\
\textbf{Taşıma altyapısı:} & CAN-FD ana omurga \\
\textbf{Durum:} & ICD Sürüm 0.3 -- v1.1 wire-format revizyonu \\
\end{tabular}
\end{center}


\section{Amaç ve Kapsam}

Bu doküman, DEOS araç içi CAN-FD haberleşme hattında kullanılan ortak adresleme,
29-bit Extended CAN Identifier yapısı, mesaj sınıfları, servisler, komut namespace'leri,
ortak veri alanı, veri tipleri ve temel mesaj formatlarını tanımlar.

Dokümanın kapsamı CAN hattında gözlenebilen protokol davranışlarıdır. Node ID, Priority,
Message Class, Service, Source, Destination, Command, Protocol Version, Sequence, Length ve
Command Payload alanlarının wire-format karşılıkları normatif olarak tanımlanır.

Gömülü yazılım mimarisi ve CAN-FD hattı dışındaki taşıma ortamlarının çalışma ayrıntıları
bu ICD'nin kapsamı dışındadır.

Bir Node ID fiziksel bir mikrodenetleyiciye birebir bağlı olmak zorunda değildir. Aynı
fiziksel cihaz üzerinde birden fazla mantıksal DEOS Node ID bulunabilir; ancak CAN hattında
her Node ID bağımsız bir protokol kimliği olarak değerlendirilir.

\section{Namespace Hiyerarşisi}

DEOS protokolünde bir mesajın kimliği tek bir ``mesaj numarası'' ile değil, birden fazla
alanın birleşimiyle belirlenir:

\begin{verbatim}
DEOS_CAN_PROTOCOL
|
+-- NODE
|   +-- SOURCE_NODE
|   +-- DESTINATION_NODE
|
+-- MESSAGE_CLASS
|
+-- SERVICE
|   +-- COMMAND
|       +-- PAYLOAD
|           +-- PARAMETER_ID   (gerektiğinde)
|           +-- VALUE          (gerektiğinde)
|
+-- PROTOCOL_METADATA
    +-- VERSION
    +-- SEQUENCE_RESERVED
    +-- LENGTH
\end{verbatim}

Buna göre bir mesajın semantik kimliği genel olarak şu bileşenlerden oluşur:

\begin{center}
\code{Priority + Message Class + Service + Destination + Source + Command}
\end{center}

Parametre yönetimi yapılan mesajlarda buna ek olarak \code{Parameter ID} kullanılır.

\section{CAN-FD Extended Identifier Yapısı}

DEOS protokolü 29-bit Extended CAN Identifier kullanır. Identifier alanı aşağıdaki
şekilde bölünür:

\begin{center}
\begin{tabular}{c c c c c}
\toprule
Bitler & Alan & Uzunluk & Açıklama \\
\midrule
28--26 & Priority & 3 bit & Arbitration önceliği \\
25--22 & Message Class & 4 bit & Mesaj davranış sınıfı \\
21--16 & Service ID & 6 bit & Fonksiyonel servis alanı \\
15--8  & Destination Node & 8 bit & Hedef düğüm adresi \\
7--0   & Source Node & 8 bit & Kaynak düğüm adresi \\
\bottomrule
\end{tabular}
\end{center}

Toplam uzunluk:

\[
3 + 4 + 6 + 8 + 8 = 29 \text{ bit}
\]

Identifier sayısal olarak aşağıdaki bağıntı ile ifade edilir:

\[
\text{CAN\_ID} =
(P \ll 26)
\;|\;
(C \ll 22)
\;|\;
(S \ll 16)
\;|\;
(D \ll 8)
\;|\;
SRC
\]

Burada:
\begin{itemize}[nosep]
    \item \(P\): Priority,
    \item \(C\): Message Class,
    \item \(S\): Service ID,
    \item \(D\): Destination Node,
    \item \(SRC\): Source Node değeridir.
\end{itemize}

CAN arbitration mantığı nedeniyle sayısal olarak daha düşük identifier daha yüksek
önceliğe sahiptir. Bu nedenle Priority alanında düşük sayısal değer daha yüksek önceliği
temsil eder.

\subsection{Önerilen Priority Registry}

\begin{center}
\begin{tabularx}{\textwidth}{c l X}
\toprule
Değer & Symbol & Kullanım \\
\midrule
0 & \code{DEOS\_PRIO\_EMERGENCY} & Acil durdurma ve en kritik safety mesajları \\
1 & \code{DEOS\_PRIO\_CRITICAL\_CONTROL} & Kritik gerçek zamanlı kontrol \\
2 & \code{DEOS\_PRIO\_CONTROL} & Normal sürüş kontrol mesajları \\
3 & \code{DEOS\_PRIO\_STATUS} & Durum ve geri besleme mesajları \\
4 & \code{DEOS\_PRIO\_CONFIG} & Konfigürasyon ve kalibrasyon \\
5 & \code{DEOS\_PRIO\_NETWORK} & Network management / heartbeat \\
6 & \code{DEOS\_PRIO\_LOW} & Düşük öncelikli servis trafiği \\
7 & \code{DEOS\_PRIO\_BACKGROUND} & Arka plan / geliştirme trafiği \\
\bottomrule
\end{tabularx}
\end{center}


\section{Node Address Registry}

Her fiziksel veya sanal DEOS haberleşme düğümüne sistem genelinde benzersiz bir 8-bit
Logical Node ID atanır. Node ID, düğümün fiziksel donanım adresi veya transport adresi
değildir; protokol içerisindeki mantıksal kimliğidir.

\subsection{Adres Aralıkları}

\begin{center}
\begin{tabularx}{\textwidth}{l X}
\toprule
Adres Aralığı & Kullanım \\
\midrule
\hex{00} & Invalid / Unassigned \\
\hex{01}--\hex{0F} & Merkezi sistem ve virtual node'lar \\
\hex{10}--\hex{1F} & Aktüatör / motor ECU düğümleri \\
\hex{20}--\hex{2F} & Güç sistemi / BMS düğümleri \\
\hex{30}--\hex{3F} & Gateway / protocol bridge düğümleri \\
\hex{40}--\hex{5F} & Sensör düğümleri için ayrılmış alan \\
\hex{60}--\hex{DF} & Gelecek genişleme \\
\hex{E0}--\hex{EF} & Geliştirme / test düğümleri \\
\hex{F0}--\hex{FE} & Diagnostik ve servis araçları \\
\hex{FF} & Broadcast \\
\bottomrule
\end{tabularx}
\end{center}

\subsection{Virtual Node Atamaları}

Virtual node, bağımsız bir fiziksel mikrodenetleyici veya bilgisayar olmak zorunda olmayan,
ancak DEOS protokolünde kendi Source ve Destination Node ID'si ile adreslenebilen mantıksal
haberleşme ucudur.

\begin{center}
\begin{tabularx}{\textwidth}{c l l X}
\toprule
Node ID & Symbol & Çalıştığı Ortam & Görev \\
\midrule
\hex{01} & \code{DEOS\_NODE\_TEXTUAL} & Raspberry Pi / Textual & Operatör arayüzü; config, command, request/response, status ve log sink \\
\hex{02} & \code{DEOS\_NODE\_GROUND\_CONTROL} & Uzak yer kontrol uygulaması & Yer kontrol haberleşme ucu \\
\hex{04} & \code{DEOS\_NODE\_MICRO\_ROS} & Raspberry Pi / micro-ROS & ROS arayüzü, telemetry ve control entegrasyonu \\
\bottomrule
\end{tabularx}
\end{center}

Textual Node yalnızca konfigürasyon amacıyla kullanılan bir node değildir. Protokolün izin
verdiği servislerde komut gönderebilir, cevap alabilir, durum bilgilerini işleyebilir ve
log mesajlarının doğrudan hedefi olabilir. Benzer şekilde micro-ROS ve Ground Control
Node'ları da ihtiyaç duydukları servisleri kaynak veya hedef olarak kullanabilir.

\subsection{Physical Node Atamaları}

\begin{center}
\begin{tabularx}{\textwidth}{c l l X}
\toprule
Node ID & Symbol & Fiziksel Cihaz & Görev \\
\midrule
\hex{03} & \code{DEOS\_NODE\_MAIN\_STM32} & STM32H563ZI & Ana fiziksel kontrol düğümü \\
\hex{10} & \code{DEOS\_NODE\_TRACTION} & STM32 Traction ECU & Tahrik sistemi kontrolü \\
\hex{11} & \code{DEOS\_NODE\_STEERING} & STM32 Steering ECU & Direksiyon sistemi kontrolü \\
\hex{12} & \code{DEOS\_NODE\_BRAKE} & STM32 Brake ECU & Fren sistemi kontrolü \\
\hex{20} & \code{DEOS\_NODE\_BMS\_MAIN} & DALY BMS \#1 & Ana batarya BMS \\
\hex{21} & \code{DEOS\_NODE\_BMS\_AUX} & DALY BMS \#2 & Yardımcı batarya BMS \\
\hex{30} & \code{DEOS\_NODE\_BMS\_GATEWAY} & CAN Gateway & Classic CAN--CAN-FD çevirici \\
\bottomrule
\end{tabularx}
\end{center}

\subsection{Servis ve Özel Node Atamaları}

\begin{center}
\begin{tabularx}{\textwidth}{c l l X}
\toprule
Node ID & Symbol & Tür & Görev \\
\midrule
\hex{F0} & \code{DEOS\_NODE\_DIAG\_TOOL} & Servis aracı & Diagnostik ve servis erişimi \\
\hex{FF} & \code{DEOS\_NODE\_BROADCAST} & Mantıksal & Tüm düğümlere yayın \\
\bottomrule
\end{tabularx}
\end{center}

\subsection{Node Enum Tanımı}

\begin{lstlisting}[style=cstyle]
typedef enum
{
    DEOS_NODE_INVALID        = 0x00,

    DEOS_NODE_TEXTUAL        = 0x01,
    DEOS_NODE_GROUND_CONTROL = 0x02,
    DEOS_NODE_MAIN_STM32     = 0x03,
    DEOS_NODE_MICRO_ROS      = 0x04,

    DEOS_NODE_TRACTION       = 0x10,
    DEOS_NODE_STEERING       = 0x11,
    DEOS_NODE_BRAKE          = 0x12,

    DEOS_NODE_BMS_MAIN       = 0x20,
    DEOS_NODE_BMS_AUX        = 0x21,

    DEOS_NODE_BMS_GATEWAY    = 0x30,

    DEOS_NODE_DIAG_TOOL      = 0xF0,
    DEOS_NODE_BROADCAST      = 0xFF
} deos_node_id_t;
\end{lstlisting}


\newpage

\section{Message Class Registry}

Message Class alanı mesajın hangi amaçla gönderildiğini tanımlar. Message Class,
fonksiyonel servisten bağımsızdır.

\begin{center}
\begin{tabularx}{\textwidth}{c l X}
\toprule
ID & Symbol & Açıklama \\
\midrule
\hex{0} & \code{DEOS\_CLASS\_SAFETY} & Acil / safety-critical mesaj \\
\hex{1} & \code{DEOS\_CLASS\_COMMAND} & Normal gerçek zamanlı kontrol komutu \\
\hex{2} & \code{DEOS\_CLASS\_STATUS} & Durum / geri besleme / telemetri \\
\hex{3} & \code{DEOS\_CLASS\_EVENT} & Fault veya olay bildirimi \\
\hex{4} & \code{DEOS\_CLASS\_REQUEST} & Cevap beklenen okuma, sorgulama veya işlem isteği \\
\hex{5} & \code{DEOS\_CLASS\_RESPONSE} & Bir REQUEST veya cevap gerektiren işleme verilen protokol cevabı \\
\hex{6} & \code{DEOS\_CLASS\_CONFIG} & Parametre ve konfigürasyon işlemleri \\
\hex{7} & \code{DEOS\_CLASS\_NETWORK} & Ağ yönetimi / heartbeat \\
\hex{8}--\hex{F} & Reserved & Gelecek kullanım \\
\bottomrule
\end{tabularx}
\end{center}

\begin{lstlisting}[style=cstyle]
typedef enum
{
    DEOS_CLASS_SAFETY   = 0x0,
    DEOS_CLASS_COMMAND  = 0x1,
    DEOS_CLASS_STATUS   = 0x2,
    DEOS_CLASS_EVENT    = 0x3,
    DEOS_CLASS_REQUEST  = 0x4,
    DEOS_CLASS_RESPONSE = 0x5,
    DEOS_CLASS_CONFIG   = 0x6,
    DEOS_CLASS_NETWORK  = 0x7
} deos_message_class_t;
\end{lstlisting}

\subsection{Message Class ve Service Arasındaki Fark}

Message Class ile Service farklı sorulara cevap verir:

\begin{center}
\begin{tabularx}{\textwidth}{l X X}
\toprule
Alan & Cevapladığı soru & Örnek \\
\midrule
Message Class & Mesaj node içerisinde hangi tür işlem akışına aittir? & COMMAND, CONFIG, NETWORK, RESPONSE \\
Service & Mesaj hangi fonksiyonel alana aittir? & STEERING, BRAKE, BMS, SYSTEM \\
Command & Bu fonksiyonel alan içerisinde tam olarak hangi işlem yapılacaktır? & SET\_TARGET\_ANGLE, SET\_PARAMETER, PING \\
\bottomrule
\end{tabularx}
\end{center}

Aynı Service farklı Message Class'lar altında kullanılabilir. Örneğin STEERING servisi için:

\begin{verbatim}
COMMAND / STEERING / SET_TARGET_ANGLE
CONFIG  / STEERING / SET_PARAMETER
REQUEST / STEERING / GET_STATUS
\end{verbatim}

Burada Service her üç mesajda da direksiyon fonksiyonel alanını belirtirken Message Class,
mesajların node içerisinde farklı işlem yollarına ayrılmasını sağlar.

Tersi de mümkündür: aynı Message Class farklı Service'lerde kullanılabilir:

\begin{verbatim}
COMMAND / TRACTION / SET_TARGET_SPEED
COMMAND / STEERING / SET_TARGET_ANGLE
COMMAND / BRAKE    / SET_TARGET
\end{verbatim}

Bu örneklerde Message Class tüm mesajların normal kontrol komutu olduğunu, Service ise hangi
alt sisteme ait olduklarını belirtir.

\section{Service Registry}

Service alanı mesajın hangi fonksiyonel domaine ait olduğunu gösterir.

\begin{center}
\begin{tabularx}{\textwidth}{c l X}
\toprule
Service ID & Symbol & Amaç \\
\midrule
\hex{00} & \code{DEOS\_SERVICE\_SYSTEM} & Genel node state/mode/reset bilgileri \\
\hex{01} & \code{DEOS\_SERVICE\_TRACTION} & Tahrik sistemi \\
\hex{02} & \code{DEOS\_SERVICE\_STEERING} & Direksiyon sistemi \\
\hex{03} & \code{DEOS\_SERVICE\_BRAKE} & Fren sistemi \\
\hex{04} & \code{DEOS\_SERVICE\_BMS} & Batarya ve BMS verileri \\
\hex{05} & \code{DEOS\_SERVICE\_CALIBRATION} & Kalibrasyon işlemleri \\
\hex{06} & \code{DEOS\_SERVICE\_DIAGNOSTIC} & Fault / sağlık / diagnostik \\
\hex{07}--\hex{3F} & Reserved & Yeni servisler için ayrılmış alan \\
\bottomrule
\end{tabularx}
\end{center}

\begin{lstlisting}[style=cstyle]
typedef enum
{
    DEOS_SERVICE_SYSTEM      = 0x00,
    DEOS_SERVICE_TRACTION    = 0x01,
    DEOS_SERVICE_STEERING    = 0x02,
    DEOS_SERVICE_BRAKE       = 0x03,
    DEOS_SERVICE_BMS         = 0x04,
    DEOS_SERVICE_CALIBRATION = 0x05,
    DEOS_SERVICE_DIAGNOSTIC  = 0x06
} deos_service_id_t;
\end{lstlisting}


\section{Node--Service İlişkisi}

Node ve Service aynı kavram değildir. Node, fiziksel veya sanal haberleşme düğümünü;
Service ise mesajın ait olduğu fonksiyonel alanı tanımlar. Bir node'un belirli bir servise
mesaj gönderebilmesi, o servisin kendisi üzerinde barındırıldığı anlamına gelmez.

Aşağıdaki tablo temel service handler sahipliğini gösterir. Her adreslenebilir DEOS node'u,
ayrıca bu ICD'de tanımlanan ortak SYSTEM/NETWORK PING--PONG davranışını desteklemek zorundadır.

\begin{center}
\begin{tabularx}{\textwidth}{l *{7}{>{\centering\arraybackslash}X}}
\toprule
Node & SYS & TRAC & STEER & BRAKE & BMS & CAL & DIAG \\
\midrule
TEXTUAL & X & -- & -- & -- & -- & -- & X \\
GROUND\_CONTROL & X & -- & -- & -- & -- & -- & X \\
MAIN\_STM32 & X & X & X & X & X & X & X \\
MICRO\_ROS & X & -- & -- & -- & -- & -- & X \\
TRACTION & X & X & -- & -- & -- & X & X \\
STEERING & X & -- & X & -- & -- & X & X \\
BRAKE & X & -- & -- & X & -- & X & X \\
BMS\_MAIN & X & -- & -- & -- & X & -- & X \\
BMS\_AUX & X & -- & -- & -- & X & -- & X \\
BMS\_GATEWAY & X & -- & -- & -- & X & -- & X \\
\bottomrule
\end{tabularx}
\end{center}

Tablodaki işaretler, ilgili node'un o servise ait mesajları desteklediğini ifade eder;
mesaj gönderme yeteneğini sınırlamaz. Örneğin \code{DEOS\_NODE\_TEXTUAL}, hedefi
\code{DEOS\_NODE\_STEERING} olan ve \code{DEOS\_SERVICE\_STEERING} kullanan bir CONFIG
veya REQUEST mesajı gönderebilir.

Örneğin direksiyon düğümüne normal sürüş komutu gönderilirken
\code{Destination = DEOS\_NODE\_STEERING}, \code{Service = DEOS\_SERVICE\_STEERING}
kullanılır. Aynı fiziksel düğüme kalibrasyon komutu gönderilirken
\code{Destination = DEOS\_NODE\_STEERING}, \code{Service = DEOS\_SERVICE\_CALIBRATION}
kullanılır.

\section{Command Namespace}

Command ID 8-bit olarak tanımlanır. Command ID tek başına global değildir;
\code{Service ID + Command ID} birleşimi komutu benzersiz hale getirir.

Genel aralık politikası:

\begin{center}
\begin{tabularx}{\textwidth}{l X}
\toprule
Command Aralığı & Kullanım \\
\midrule
\hex{00} & Invalid / Reserved \\
\hex{01}--\hex{7F} & Service-specific komutlar \\
\hex{80}--\hex{DF} & Gelecek service-specific genişleme \\
\hex{E0}--\hex{E4} & Ortak Parameter Registry komutları \\
\hex{E5}--\hex{EF} & Ortak yönetim komutları için reserved \\
\hex{F0}--\hex{FE} & Geliştirme / deneysel \\
\hex{FF} & Reserved \\
\bottomrule
\end{tabularx}
\end{center}

\subsection{SYSTEM Service Komutları}

\begin{center}
\begin{tabular}{c l}
\toprule
ID & Symbol \\
\midrule
\hex{01} & \code{DEOS\_CMD\_SYSTEM\_GET\_STATE} \\
\hex{02} & \code{DEOS\_CMD\_SYSTEM\_SET\_STATE} \\
\hex{03} & \code{DEOS\_CMD\_SYSTEM\_GET\_MODE} \\
\hex{04} & \code{DEOS\_CMD\_SYSTEM\_SET\_MODE} \\
\hex{05} & \code{DEOS\_CMD\_SYSTEM\_RESET\_NODE} \\
\hex{06} & \code{DEOS\_CMD\_SYSTEM\_GET\_NODE\_INFO} \\
\hex{07} & \code{DEOS\_CMD\_SYSTEM\_HEARTBEAT} \\
\hex{08} & \code{DEOS\_CMD\_SYSTEM\_PING} \\
\hex{09} & \code{DEOS\_CMD\_SYSTEM\_PONG} \\
\bottomrule
\end{tabular}
\end{center}

\begin{lstlisting}[style=cstyle]
typedef enum
{
    DEOS_CMD_SYSTEM_GET_STATE     = 0x01,
    DEOS_CMD_SYSTEM_SET_STATE     = 0x02,
    DEOS_CMD_SYSTEM_GET_MODE      = 0x03,
    DEOS_CMD_SYSTEM_SET_MODE      = 0x04,
    DEOS_CMD_SYSTEM_RESET_NODE    = 0x05,
    DEOS_CMD_SYSTEM_GET_NODE_INFO = 0x06,
    DEOS_CMD_SYSTEM_HEARTBEAT     = 0x07,
    DEOS_CMD_SYSTEM_PING          = 0x08,
    DEOS_CMD_SYSTEM_PONG          = 0x09
} deos_system_command_t;
\end{lstlisting}

\subsection{TRACTION Service Komutları}

\begin{center}
\begin{tabular}{c l}
\toprule
ID & Symbol \\
\midrule
\hex{01} & \code{DEOS\_CMD\_TRACTION\_SET\_TARGET\_SPEED} \\
\hex{02} & \code{DEOS\_CMD\_TRACTION\_GET\_STATUS} \\
\bottomrule
\end{tabular}
\end{center}

\begin{lstlisting}[style=cstyle]
typedef enum
{
    DEOS_CMD_TRACTION_SET_TARGET_SPEED = 0x01,
    DEOS_CMD_TRACTION_GET_STATUS       = 0x02
} deos_traction_command_t;
\end{lstlisting}

\subsection{STEERING Service Komutları}

\begin{center}
\begin{tabular}{c l}
\toprule
ID & Symbol \\
\midrule
\hex{01} & \code{DEOS\_CMD\_STEERING\_SET\_TARGET\_ANGLE} \\
\hex{02} & \code{DEOS\_CMD\_STEERING\_GET\_STATUS} \\
\bottomrule
\end{tabular}
\end{center}

\begin{lstlisting}[style=cstyle]
typedef enum
{
    DEOS_CMD_STEERING_SET_TARGET_ANGLE = 0x01,
    DEOS_CMD_STEERING_GET_STATUS       = 0x02
} deos_steering_command_t;
\end{lstlisting}

\subsection{BRAKE Service Komutları}

\begin{center}
\begin{tabular}{c l}
\toprule
ID & Symbol \\
\midrule
\hex{01} & \code{DEOS\_CMD\_BRAKE\_SET\_TARGET} \\
\hex{02} & \code{DEOS\_CMD\_BRAKE\_GET\_STATUS} \\
\bottomrule
\end{tabular}
\end{center}

\begin{lstlisting}[style=cstyle]
typedef enum
{
    DEOS_CMD_BRAKE_SET_TARGET = 0x01,
    DEOS_CMD_BRAKE_GET_STATUS = 0x02
} deos_brake_command_t;
\end{lstlisting}

\subsection{BMS Service Komutları}

\begin{center}
\begin{tabular}{c l}
\toprule
ID & Symbol \\
\midrule
\hex{01} & \code{DEOS\_CMD\_BMS\_GET\_STATUS} \\
\hex{02} & \code{DEOS\_CMD\_BMS\_STATUS} \\
\bottomrule
\end{tabular}
\end{center}

\begin{lstlisting}[style=cstyle]
typedef enum
{
    DEOS_CMD_BMS_GET_STATUS = 0x01,
    DEOS_CMD_BMS_STATUS     = 0x02
} deos_bms_command_t;
\end{lstlisting}

\subsection{CALIBRATION Service Komutları}

\begin{center}
\begin{tabular}{c l}
\toprule
ID & Symbol \\
\midrule
\hex{01} & \code{DEOS\_CMD\_CAL\_START} \\
\hex{02} & \code{DEOS\_CMD\_CAL\_ABORT} \\
\hex{03} & \code{DEOS\_CMD\_CAL\_GET\_STATUS} \\
\hex{04} & \code{DEOS\_CMD\_CAL\_SAVE\_RESULT} \\
\hex{05} & \code{DEOS\_CMD\_CAL\_CLEAR\_RESULT} \\
\bottomrule
\end{tabular}
\end{center}

\begin{lstlisting}[style=cstyle]
typedef enum
{
    DEOS_CMD_CAL_START        = 0x01,
    DEOS_CMD_CAL_ABORT        = 0x02,
    DEOS_CMD_CAL_GET_STATUS   = 0x03,
    DEOS_CMD_CAL_SAVE_RESULT  = 0x04,
    DEOS_CMD_CAL_CLEAR_RESULT = 0x05
} deos_calibration_command_t;
\end{lstlisting}

\subsection{DIAGNOSTIC Service Komutları}

\begin{center}
\begin{tabular}{c l}
\toprule
ID & Symbol \\
\midrule
\hex{01} & \code{DEOS\_CMD\_DIAG\_GET\_HEALTH} \\
\hex{02} & \code{DEOS\_CMD\_DIAG\_GET\_FAULTS} \\
\hex{03} & \code{DEOS\_CMD\_DIAG\_CLEAR\_FAULTS} \\
\bottomrule
\end{tabular}
\end{center}

\begin{lstlisting}[style=cstyle]
typedef enum
{
    DEOS_CMD_DIAG_GET_HEALTH   = 0x01,
    DEOS_CMD_DIAG_GET_FAULTS   = 0x02,
    DEOS_CMD_DIAG_CLEAR_FAULTS = 0x03
} deos_diagnostic_command_t;
\end{lstlisting}

\section{CONFIG Message Class ve Parameter Registry}

Konfigürasyon işlemleri için ayrı bir CONFIG servisi kullanılmaz. Bunun yerine
\code{Message Class = CONFIG} seçilir ve hedef fonksiyonel servis korunur. REQUEST ve RESPONSE
sınıfları ise protokolde ayrıca bulunmaya devam eder; CONFIG bunların yerine geçen bir sınıf değildir.

Örneğin direksiyon PID katsayısı değiştirilirken:

\begin{verbatim}
Destination   = DEOS_NODE_STEERING
Message Class = DEOS_CLASS_CONFIG
Service       = DEOS_SERVICE_STEERING
Command       = DEOS_CMD_SET_PARAMETER
Parameter ID  = DEOS_PARAM_STEERING_PID_KP
\end{verbatim}

Bu yapı sayesinde konfigürasyon işlemi hangi alt sisteme ait olduğunu kaybetmez.

\subsection{Ortak Parameter Command ID'leri}

Aşağıdaki command ID'leri configurable servislerde ortak olarak kullanılır:

\begin{center}
\begin{tabular}{c l}
\toprule
ID & Symbol \\
\midrule
\hex{E0} & \code{DEOS\_CMD\_GET\_PARAMETER} \\
\hex{E1} & \code{DEOS\_CMD\_SET\_PARAMETER} \\
\hex{E2} & \code{DEOS\_CMD\_SAVE\_PARAMETERS} \\
\hex{E3} & \code{DEOS\_CMD\_RESTORE\_DEFAULTS} \\
\hex{E4} & \code{DEOS\_CMD\_GET\_PARAMETER\_INFO} \\
\bottomrule
\end{tabular}
\end{center}

\begin{lstlisting}[style=cstyle]
typedef enum
{
    DEOS_CMD_GET_PARAMETER       = 0xE0,
    DEOS_CMD_SET_PARAMETER       = 0xE1,
    DEOS_CMD_SAVE_PARAMETERS     = 0xE2,
    DEOS_CMD_RESTORE_DEFAULTS    = 0xE3,
    DEOS_CMD_GET_PARAMETER_INFO  = 0xE4
} deos_parameter_command_t;
\end{lstlisting}

\subsection{Parameter ID Alanı}

Parameter ID 16-bit olarak tanımlanır. Parameter ID global değildir;
\code{Service ID + Parameter ID} birlikte benzersizdir.

Örnek Steering Parameter Registry:

\begin{center}
\begin{tabularx}{\textwidth}{c l X}
\toprule
ID & Symbol & Açıklama \\
\midrule
\hex{0001} & \code{DEOS\_PARAM\_STEERING\_PID\_KP} & PID oransal katsayısı \\
\hex{0002} & \code{DEOS\_PARAM\_STEERING\_PID\_KI} & PID integral katsayısı \\
\hex{0003} & \code{DEOS\_PARAM\_STEERING\_PID\_KD} & PID türev katsayısı \\
\hex{0010} & \code{DEOS\_PARAM\_STEERING\_ZERO\_OFFSET} & Direksiyon sıfır ofseti \\
\hex{0011} & \code{DEOS\_PARAM\_STEERING\_MIN\_ANGLE} & Minimum izinli açı \\
\hex{0012} & \code{DEOS\_PARAM\_STEERING\_MAX\_ANGLE} & Maksimum izinli açı \\
\hex{0013} & \code{DEOS\_PARAM\_STEERING\_CURRENT\_LIMIT} & Akım limiti \\
\bottomrule
\end{tabularx}
\end{center}

Örnek Traction Parameter Registry:

\begin{center}
\begin{tabularx}{\textwidth}{c l X}
\toprule
ID & Symbol & Açıklama \\
\midrule
\hex{0001} & \code{DEOS\_PARAM\_TRACTION\_PID\_KP} & PID oransal katsayısı \\
\hex{0002} & \code{DEOS\_PARAM\_TRACTION\_PID\_KI} & PID integral katsayısı \\
\hex{0003} & \code{DEOS\_PARAM\_TRACTION\_PID\_KD} & PID türev katsayısı \\
\hex{0010} & \code{DEOS\_PARAM\_TRACTION\_MAX\_CURRENT} & Maksimum akım \\
\hex{0011} & \code{DEOS\_PARAM\_TRACTION\_MAX\_SPEED} & Maksimum hız \\
\hex{0012} & \code{DEOS\_PARAM\_TRACTION\_ACCEL\_LIMIT} & İvme limiti \\
\bottomrule
\end{tabularx}
\end{center}

Örnek Brake Parameter Registry:

\begin{center}
\begin{tabularx}{\textwidth}{c l X}
\toprule
ID & Symbol & Açıklama \\
\midrule
\hex{0001} & \code{DEOS\_PARAM\_BRAKE\_PID\_KP} & Gelecek kullanım / gerekli olması halinde \\
\hex{0002} & \code{DEOS\_PARAM\_BRAKE\_PID\_KI} & Gelecek kullanım / gerekli olması halinde \\
\hex{0003} & \code{DEOS\_PARAM\_BRAKE\_PID\_KD} & Gelecek kullanım / gerekli olması halinde \\
\hex{0010} & \code{DEOS\_PARAM\_BRAKE\_MIN\_POSITION} & Minimum aktüatör konumu \\
\hex{0011} & \code{DEOS\_PARAM\_BRAKE\_MAX\_POSITION} & Maksimum aktüatör konumu \\
\hex{0012} & \code{DEOS\_PARAM\_BRAKE\_CURRENT\_LIMIT} & Akım limiti \\
\bottomrule
\end{tabularx}
\end{center}

\section{State ve Mode Registry}

State, düğümün anlık çalışma durumunu; Mode ise düğümün hangi çalışma davranışı altında
olduğunu ifade eder. İki alan birbirinden bağımsızdır.

\subsection{State Enum}

\begin{center}
\begin{tabular}{c l}
\toprule
Değer & Symbol \\
\midrule
\hex{00} & \code{DEOS\_STATE\_INIT} \\
\hex{01} & \code{DEOS\_STATE\_STANDBY} \\
\hex{02} & \code{DEOS\_STATE\_READY} \\
\hex{03} & \code{DEOS\_STATE\_ACTIVE} \\
\hex{04} & \code{DEOS\_STATE\_CALIBRATING} \\
\hex{05} & \code{DEOS\_STATE\_SAFE} \\
\hex{06} & \code{DEOS\_STATE\_FAULT} \\
\hex{FF} & \code{DEOS\_STATE\_UNKNOWN} \\
\bottomrule
\end{tabular}
\end{center}

\begin{lstlisting}[style=cstyle]
typedef enum
{
    DEOS_STATE_INIT        = 0x00,
    DEOS_STATE_STANDBY     = 0x01,
    DEOS_STATE_READY       = 0x02,
    DEOS_STATE_ACTIVE      = 0x03,
    DEOS_STATE_CALIBRATING = 0x04,
    DEOS_STATE_SAFE        = 0x05,
    DEOS_STATE_FAULT       = 0x06,
    DEOS_STATE_UNKNOWN     = 0xFF
} deos_state_t;
\end{lstlisting}

\subsection{Mode Enum}

\begin{center}
\begin{tabular}{c l}
\toprule
Değer & Symbol \\
\midrule
\hex{00} & \code{DEOS\_MODE\_NORMAL} \\
\hex{01} & \code{DEOS\_MODE\_MANUAL} \\
\hex{02} & \code{DEOS\_MODE\_AUTONOMOUS} \\
\hex{03} & \code{DEOS\_MODE\_SERVICE} \\
\hex{04} & \code{DEOS\_MODE\_TEST} \\
\hex{FF} & \code{DEOS\_MODE\_UNKNOWN} \\
\bottomrule
\end{tabular}
\end{center}

\begin{lstlisting}[style=cstyle]
typedef enum
{
    DEOS_MODE_NORMAL     = 0x00,
    DEOS_MODE_MANUAL     = 0x01,
    DEOS_MODE_AUTONOMOUS = 0x02,
    DEOS_MODE_SERVICE    = 0x03,
    DEOS_MODE_TEST       = 0x04,
    DEOS_MODE_UNKNOWN    = 0xFF
} deos_mode_t;
\end{lstlisting}

\section{Uygulama Veri Alanı}

DEOS CAN-FD data alanında Source, Destination, Service ve Priority tekrar taşınmaz; bu
alanlar 29-bit CAN Identifier içerisinde bulunur. Uygulama seviyesinde ek senkronizasyon
kelimesi veya ek CRC tanımlanmaz.

DEOS v1.1 ortak veri alanı aşağıdaki gibidir:

\begin{center}
\begin{tabular}{c c c c}
\toprule
Byte & Alan & Uzunluk & Açıklama \\
\midrule
0 & Protocol Version & 1 byte & Protokol major/minor sürümü \\
1 & Sequence / Reserved & 1 byte & Source Node başına TX sequence değeri \\
2 & Command ID & 1 byte & İlgili Service namespace'i içindeki komut \\
3 & Length & 1 byte & Geçerli command-specific payload byte sayısı \\
4...N & Command Payload & Değişken & Komuta özgü veri \\
\bottomrule
\end{tabular}
\end{center}

\subsection{Length Alanı}

Length alanı \code{uint8\_t} olarak tanımlanır ve yalnızca command-specific payload'un
geçerli byte sayısını belirtir. Version, Sequence, Command ve Length alanlarının kendisi
bu değere dahil değildir. Geçerli aralık \code{0..60} byte'tır.

CAN-FD DLC yuvarlaması nedeniyle oluşan padding byte'ları Length değerine dahil değildir
ve command payload olarak yorumlanmaz.

Maksimum CAN-FD data alanı 64 byte ve ortak DEOS header'ı 4 byte olduğundan maksimum
command-specific payload uzunluğu 60 byte'tır.

\subsection{Protocol Version}

Version byte iki adet 4-bit alana ayrılır:

\begin{center}
\begin{tabular}{c c}
\toprule
Bitler & Anlam \\
\midrule
7--4 & Major Version \\
3--0 & Minor Version \\
\bottomrule
\end{tabular}
\end{center}

Aktif wire-format sürümü:

\[
\text{DEOS Protocol v1.1} = \hex{11}
\]

\hex{10} değeri v1.0 eski wire-formatını, \hex{11} değeri ise Length alanını içeren v1.1
wire-formatını ifade eder. Bu iki wire-format birbiriyle uyumlu değildir. v1.1 formatında
geçerli bir frame'in Protocol Version alanı \hex{11} olmalıdır.

\subsection{Sequence Alanı}

Sequence alanı 8-bit'tir. Her Source Node kimliği kendi transmit sequence değerini kullanır.
Sayaç \hex{FF} sonrasında \hex{00} değerine dönebilir.

v1.1'de Sequence alanı protokol geçerliliği açısından bilgilendirici/reserved semantiğini
korur. Receiver, sequence atlamasını veya tekrarını zorunlu olarak paket kaybı, duplicate
frame veya protokol hatası olarak yorumlamak zorunda değildir.

\section{Veri Tipleri ve Byte Sırası}

DEOS protokolünde çok baytlı alanlar için varsayılan byte sırası \textbf{Little-Endian}
olarak tanımlanır.

Wire formatında temel tercih, doğrudan IEEE-754 floating-point taşımak yerine ölçekli
integer değer kullanmaktır. Böylece bir sinyalin fiziksel karşılığı açık şekilde
\code{raw value + scale + offset + unit} ile tanımlanır.

Genel fiziksel değer dönüşümü:

\[
PhysicalValue = RawValue \times Scale + Offset
\]

\subsection{Normatif Veri Tipleri}

\begin{center}
\begin{tabularx}{\textwidth}{l c X}
\toprule
Tip & Boyut & Kullanım \\
\midrule
\code{uint8\_t}  & 1 byte & Enum, yüzde, küçük unsigned değerler \\
\code{int8\_t}   & 1 byte & Küçük signed değerler \\
\code{uint16\_t} & 2 byte & Ölçekli pozitif fiziksel değerler \\
\code{int16\_t}  & 2 byte & Ölçekli signed fiziksel değerler \\
\code{uint32\_t} & 4 byte & Sayaç / geniş pozitif değerler \\
\code{int32\_t}  & 4 byte & Geniş signed fiziksel değerler \\
\code{float32}   & 4 byte & Yalnız ilgili mesaj tanımında açıkça belirtilirse IEEE-754 \\
\bottomrule
\end{tabularx}
\end{center}

Floating-point kullanımına genel yasak konulmamıştır; ancak bir alan \code{float32}
olarak tanımlanmamışsa integer tabanlı ölçekleme varsayılır.

\subsection{Başlangıç Fiziksel Sinyal Formatları}

\begin{center}
\begin{tabularx}{\textwidth}{l l l l X}
\toprule
Sinyal & Wire Type & Scale & Unit & Örnek \\
\midrule
Steering target angle & \code{int16\_t} & 0.01 & degree & 1250 = 12.50 deg \\
Vehicle target speed & \code{uint16\_t} & 0.01 & m/s & 725 = 7.25 m/s \\
Brake target & \code{uint16\_t} & 0.01 & \% & 7500 = 75.00\% \\
BMS voltage & \code{uint16\_t} & 0.01 & V & 7425 = 74.25 V \\
BMS current & \code{int16\_t} & 0.1 & A & -125 = -12.5 A \\
BMS SOC & \code{uint16\_t} & 0.01 & \% & 8530 = 85.30\% \\
Temperature & \code{int16\_t} & 0.1 & degC & 253 = 25.3 degC \\
\bottomrule
\end{tabularx}
\end{center}

Bu tablo başlangıç wire-format standardını tanımlar. Nihai min/max sınırlar ve
invalid değerler ilgili komut veya sinyal tanımında ayrıca belirtilmelidir.

\section{CAN-FD DLC, Payload Uzunluğu ve BRS}

CAN-FD DLC alanı, frame'in fiziksel CAN-FD data field boyutunu belirtir. DEOS v1.1
Length alanı ise yalnızca geçerli command-specific payload byte sayısını belirtir. Bu iki
alan aynı anlamda değildir.

8 byte üzerindeki CAN-FD data alanları 12, 16, 20, 24, 32, 48 veya 64 byte fiziksel
boyutlara yuvarlanabilir. Bu nedenle DEOS Length değeri fiziksel CAN-FD data field
boyutundan bağımsız olarak taşınır.

Örnek olarak 10 byte command payload için:

\begin{verbatim}
DEOS common header  = 4 byte
Command payload     = 10 byte
Semantic data size  = 14 byte
CAN-FD data size    = 16 byte
Padding             = 2 byte
\end{verbatim}

DLC yuvarlaması nedeniyle semantic data sonrasında kalan kullanılmayan byte'lar \hex{00}
olarak gönderilir ve alıcı tarafından payload'ın parçası olarak değerlendirilmez.

DEOS CAN-FD frame'leri 29-bit Extended Identifier, CAN-FD formatı ve Bit Rate Switching
(BRS) kullanır. Nominal/arbitration bitrate ile data-phase bitrate değerleri deployment
parametreleridir ve bu ICD revizyonunda sabitlenmemiştir.

\section{REQUEST--RESPONSE ve CONFIG Ayrımı}

REQUEST, RESPONSE ve CONFIG birbirinden bağımsız mesaj sınıflarıdır ve farklı amaçlara hizmet eder.

\begin{itemize}
    \item \code{REQUEST}: Bir düğümden veri, durum veya belirli bir işlem sonucu talep etmek için kullanılır.
    \item \code{RESPONSE}: REQUEST sınıfındaki bir mesaja veya cevap gerektiren başka bir işleme karşılık dönen cevabı taşır.
    \item \code{CONFIG}: Bir düğümün çalışma parametrelerini, ayarlarını veya kalıcı yapılandırmasını değiştirmeye yönelik mesajları taşır.
\end{itemize}

Örneğin bir parametreyi okumak için REQUEST, aynı parametreyi değiştirmek için CONFIG
kullanılabilir:

\begin{verbatim}
REQUEST / STEERING / GET_PARAMETER / PID_KP
CONFIG  / STEERING / SET_PARAMETER / PID_KP
\end{verbatim}

Her iki işlem için de gerektiğinde RESPONSE sınıfı ile sonuç döndürülebilir. Bu ayrım,
mesajın ``cevap bekleyip beklemediği'' ile ``konfigürasyon amacı taşıyıp taşımadığı''
bilgilerinin birbirine karıştırılmasını önler.

\section{Response Altyapısı}

RESPONSE Message Class protokolde ayrılmıştır; ancak hangi komutların zorunlu response
üreteceği bu registry dokümanında kesinleştirilmemiştir.

Özellikle aşağıdaki işlemler gelecekte response kullanılabilecek adaylardır:
\begin{itemize}[nosep]
    \item Parameter GET/SET,
    \item state veya mode değişimi,
    \item kalibrasyon başlatma / kaydetme,
    \item fault reset / clear,
    \item node reset,
    \item servis tipi tek-seferlik yönetim komutları.
\end{itemize}

\subsection{Generic Response Payload}

\code{DEOS\_CLASS\_RESPONSE} mesajları varsayılan olarak standart bir veri yerleşimi kullanır. Eğer ilgili komut özel bir ek veri döndürmüyorsa, yanıt mesajının içeriği her zaman şu formatta olur:

\begin{center}
\begin{tabularx}{\textwidth}{c l c X}
\toprule
Byte & Alan & Tip & Açıklama \\
\midrule
0 & Result Code & \code{uint8\_t} & İşlem sonucu (\code{deos\_result\_t}) \\
1..N & Ek Veri (Opsiyonel) & Değişken & Komuta özgü veri \\
\bottomrule
\end{tabularx}
\end{center}

Result code \hex{00} (SUCCESS) olduğunda işlem başarılı sayılır. Bu genel yapı sayesinde, her komut için özel bir response formatı tasarlamaya gerek kalmadan hata yönetimi standartlaştırılmıştır.

Aşağıda standart response result kodları listelenmiştir:

\begin{center}
\begin{tabular}{c l}
\toprule
Kod & Anlam \\
\midrule
\hex{00} & SUCCESS \\
\hex{01} & UNSUPPORTED \\
\hex{02} & INVALID\_PARAMETER \\
\hex{03} & OUT\_OF\_RANGE \\
\hex{04} & NOT\_ALLOWED \\
\hex{05} & BUSY \\
\hex{06} & NOT\_READY \\
\hex{07} & INTERNAL\_ERROR \\
\bottomrule
\end{tabular}
\end{center}

Bu alanın kesin payload yerleşimi sonraki ICD revizyonunda mesaj bazında ayrıntılandırılacaktır.


\section{PING--PONG Ağ Erişilebilirlik Mekanizması}

DEOS protokolünde tüm adreslenebilir physical ve virtual node'lar ortak bir PING--PONG
mekanizması uygular. Bir node, erişmek istediği başka bir node'a unicast PING gönderebilir.
Geçerli PING mesajını alan hedef node, aynı Ping ID değerini içeren PONG mesajını kaynak
node'a geri göndermek zorundadır.

PING ve PONG uygulama REQUEST/RESPONSE mekanizmasından ayrı bir network-management
mekanizmasıdır. Her iki mesaj da aşağıdaki alanları kullanır:

\begin{verbatim}
Priority      = DEOS_PRIO_NETWORK
Message Class = DEOS_CLASS_NETWORK
Service       = DEOS_SERVICE_SYSTEM
\end{verbatim}

PING için \code{Command = DEOS\_CMD\_SYSTEM\_PING}; PONG için
\code{Command = DEOS\_CMD\_SYSTEM\_PONG} kullanılır.

\subsection{PING Mesaj Formatı}

PING mesajının uygulama veri alanı aşağıdaki gibi tanımlanır:

\begin{center}
\begin{tabularx}{\textwidth}{c l c X}
\toprule
Byte & Alan & Tip & Açıklama \\
\midrule
0 & Protocol Version & \code{uint8\_t} & DEOS v1.1 için \hex{11} \\
1 & Sequence & \code{uint8\_t} & Gönderen node'un TX sequence değeri \\
2 & Command & \code{uint8\_t} & \code{DEOS\_CMD\_SYSTEM\_PING} = \hex{08} \\
3 & Length & \code{uint8\_t} & \hex{04} \\
4--7 & Ping ID & \code{uint32\_t} & Gönderen tarafından üretilen correlation değeri \\
\bottomrule
\end{tabularx}
\end{center}

Toplam semantic CAN-FD data length 8 byte'tır. Ping ID, DEOS genel byte-order kuralına uygun olarak
Little-Endian taşınır. Gönderen node, aynı anda birden fazla PING işlemi yürütüyorsa aktif
işlemleri ayırt edecek Ping ID değerleri kullanmalıdır.

PING CAN Identifier alanları:

\begin{verbatim}
Priority      = DEOS_PRIO_NETWORK
Message Class = DEOS_CLASS_NETWORK
Service       = DEOS_SERVICE_SYSTEM
Destination   = ping atılan node
Source        = ping gönderen node
\end{verbatim}

\subsection{PONG Mesaj Formatı}

PONG mesajı, alınan PING mesajındaki 32-bit Ping ID değerini değiştirmeden geri taşır.

\begin{center}
\begin{tabularx}{\textwidth}{c l c X}
\toprule
Byte & Alan & Tip & Açıklama \\
\midrule
0 & Protocol Version & \code{uint8\_t} & DEOS v1.1 için \hex{11} \\
1 & Sequence & \code{uint8\_t} & PONG gönderen node'un kendi TX sequence değeri \\
2 & Command & \code{uint8\_t} & \code{DEOS\_CMD\_SYSTEM\_PONG} = \hex{09} \\
3 & Length & \code{uint8\_t} & \hex{04} \\
4--7 & Ping ID & \code{uint32\_t} & PING'den aynen echo edilen correlation değeri \\
\bottomrule
\end{tabularx}
\end{center}

PONG CAN Identifier alanları:

\begin{verbatim}
Priority      = DEOS_PRIO_NETWORK
Message Class = DEOS_CLASS_NETWORK
Service       = DEOS_SERVICE_SYSTEM
Destination   = PING mesajının Source Node değeri
Source        = PING'i alan hedef node
\end{verbatim}

PONG mesajındaki Sequence alanı PING mesajındaki Sequence değerini kopyalamaz. Her node
kendi transmit sequence sayacını kullanmaya devam eder. PING ile PONG arasındaki transaction
correlation yalnızca Ping ID üzerinden yapılır.

\subsection{Zorunlu Node Davranışı}

Geçerli ve kendisine yöneltilmiş bir PING alan her DEOS node'u aşağıdaki davranışı uygular:

\begin{enumerate}[nosep]
    \item PING mesajının Source Node ve Ping ID alanlarını okur.
    \item Yeni bir PONG mesajı oluşturur.
    \item PONG Source alanını kendi Node ID'si yapar.
    \item PONG Destination alanını PING Source değerine eşitler.
    \item PING içerisindeki Ping ID değerini değiştirmeden PONG payload'ına kopyalar.
    \item Kendi normal TX Sequence değerini kullanarak PONG'u gönderir.
\end{enumerate}

\subsection{RTT Ölçümü}

Round-trip time ölçümü gönderen node üzerinde yapılır. PING gönderildiği anda yerel zaman
kaydedilir; aynı Ping ID'yi taşıyan PONG alındığında iki yerel zaman arasındaki fark RTT
olarak hesaplanır:

\[
RTT = T_{pong} - T_{ping}
\]

Bu yöntem node'lar arasında saat senkronizasyonu gerektirmez. Timeout süresi transport veya
uygulama gereksinimine göre ayrıca tanımlanabilir; v0.3 kapsamında tek bir global timeout
süresi dondurulmamıştır.

\subsection{Broadcast PING Kuralı}

\code{DEOS\_CMD\_SYSTEM\_PING} mesajı v1.1 protokolünde yalnızca unicast Destination Node
ile kullanılmalıdır. \code{Destination = DEOS\_NODE\_BROADCAST} ile PING gönderimi geçersizdir.
Bu kural, çok sayıda node'un aynı anda PONG üretmesinden doğacak response burst trafiğini
önlemek için tanımlanmıştır.

Bir uygulama tüm node'ların erişilebilirliğini taramak istiyorsa her hedef Node ID için ayrı
unicast PING işlemi başlatmalıdır.

\subsection{HEARTBEAT ile PING--PONG Ayrımı}

HEARTBEAT, bir node'un periyodik veya olay tabanlı olarak kendi yaşam bilgisini yayınlamasıdır.
PING--PONG ise belirli bir node'un başka belirli bir node'a aktif erişilebilirlik sorgusu
yapmasıdır. Bu iki mekanizma birbirinin yerine kullanılmaz.

\section{Mantıksal Node Kimliği Kuralı}

CAN hattında her Source ve Destination Node ID bağımsız bir DEOS protokol kimliğidir.
Aynı fiziksel cihaz üzerinde birden fazla mantıksal Node ID bulunabilse bile, bir mesaj
hangi Node ID adına üretiliyorsa CAN Identifier içerisindeki Source alanında o Node ID
kullanılmalıdır.

Bir node başka bir node'un kimliğiyle cevap üretmemelidir. Örneğin
\code{Destination = DEOS\_NODE\_BMS\_MAIN} olan geçerli bir PING'e verilen PONG mesajında
\code{Source = DEOS\_NODE\_BMS\_MAIN} kullanılmalıdır.

Bu kural fiziksel yerleşimden bağımsızdır; CAN hattındaki diğer katılımcılar mantıksal
Node ID'leri bağımsız haberleşme uçları olarak değerlendirir.

\section{Örnek Mesaj Kimlikleri}

\subsection{Direksiyon Hedef Açısı}

\begin{verbatim}
Source        = MAIN_STM32
Destination   = STEERING
Priority      = CONTROL
Message Class = COMMAND
Service       = STEERING
Command       = SET_TARGET_ANGLE
Payload       = int16_t target_angle, scale 0.01 deg
\end{verbatim}

\subsection{Direksiyon PID Kp Güncellemesi}

\begin{verbatim}
Source        = MAIN_STM32
Destination   = STEERING
Priority      = CONFIG
Message Class = CONFIG
Service       = STEERING
Command       = SET_PARAMETER
Parameter ID  = STEERING_PID_KP
Value         = ilgili parameter taniminda belirtilen wire type
\end{verbatim}

\subsection{Direksiyon Kalibrasyonu Başlatma}

\begin{verbatim}
Source        = MAIN_STM32
Destination   = STEERING
Message Class = COMMAND
Service       = CALIBRATION
Command       = START
\end{verbatim}

\subsection{BMS Durum Mesajı}

\begin{verbatim}
Source        = BMS_MAIN
Destination   = MAIN_STM32
Message Class = STATUS
Service       = BMS
Command       = STATUS
\end{verbatim}

Gateway kullanılması bu mantıksal kimliği değiştirmez.


\subsection{Ground Control'dan Textual Node'a PING}

\begin{verbatim}
Source        = DEOS_NODE_GROUND_CONTROL
Destination   = DEOS_NODE_TEXTUAL
Priority      = DEOS_PRIO_NETWORK
Message Class = DEOS_CLASS_NETWORK
Service       = DEOS_SERVICE_SYSTEM
Command       = DEOS_CMD_SYSTEM_PING
Ping ID       = 0x00001234
\end{verbatim}

Textual Node aşağıdaki PONG mesajını üretir:

\begin{verbatim}
Source        = DEOS_NODE_TEXTUAL
Destination   = DEOS_NODE_GROUND_CONTROL
Priority      = DEOS_PRIO_NETWORK
Message Class = DEOS_CLASS_NETWORK
Service       = DEOS_SERVICE_SYSTEM
Command       = DEOS_CMD_SYSTEM_PONG
Ping ID       = 0x00001234
\end{verbatim}


\section{İsimlendirme Standardı}

Tüm protokol sabitlerinde ortak prefix yapısı kullanılır:

\begin{verbatim}
DEOS_NODE_...
DEOS_PRIO_...
DEOS_CLASS_...
DEOS_SERVICE_...
DEOS_CMD_<SERVICE>_...
DEOS_PARAM_<SERVICE>_...
DEOS_STATE_...
DEOS_MODE_...
\end{verbatim}

Örnekler:

\begin{verbatim}
DEOS_NODE_TEXTUAL
DEOS_NODE_GROUND_CONTROL
DEOS_NODE_MAIN_STM32
DEOS_NODE_MICRO_ROS
DEOS_NODE_STEERING
DEOS_CLASS_CONFIG
DEOS_SERVICE_STEERING
DEOS_CMD_STEERING_SET_TARGET_ANGLE
DEOS_CMD_SET_PARAMETER
DEOS_PARAM_STEERING_PID_KP
DEOS_STATE_ACTIVE
DEOS_MODE_AUTONOMOUS
\end{verbatim}


\section{Özet Registry}

Mevcut CAN-FD protokol kararları aşağıdaki şekilde özetlenebilir:

\begin{itemize}
    \item Ana araç içi haberleşme hattı CAN-FD'dir.
    \item 29-bit Extended CAN Identifier kullanılır.
    \item CAN-FD frame'lerinde BRS kullanılır.
    \item Priority, Message Class, Service, Destination ve Source alanları CAN Identifier içerisinde taşınır.
    \item Node ID, sistem genelinde kullanılan 8-bit mantıksal protokol kimliğidir.
    \item Aynı fiziksel cihaz birden fazla mantıksal Node ID barındırabilir; CAN hattında her biri bağımsız kimliktir.
    \item DEOS Protocol v1.1 sürüm değeri \hex{11}'dir.
    \item Ortak DEOS data header'ı Version, Sequence, Command ve Length olmak üzere 4 byte'tır.
    \item Command-specific maksimum payload uzunluğu 60 byte'tır.
    \item Length, yalnızca geçerli command-specific payload byte sayısını belirtir.
    \item DLC fiziksel CAN-FD data field boyutunu; Length ise semantic DEOS payload boyutunu ifade eder.
    \item DLC yuvarlaması sonucu oluşan kullanılmayan byte'lar \hex{00} padding olarak gönderilir.
    \item Uygulama seviyesinde ek CRC veya ayrıca sync header kullanılmaz.
    \item Sequence alanı Source Node başına kullanılır; v1.1'de gap/duplicate kontrolü zorunlu değildir.
    \item REQUEST, RESPONSE ve CONFIG birbirinden ayrı Message Class değerleridir.
    \item CONFIG ayrı bir Service değildir.
    \item Parameter ID 16-bit'tir ve Service namespace'i içinde anlam kazanır.
    \item Varsayılan çok baytlı wire byte sırası Little-Endian'dır.
    \item Fiziksel değerlerde varsayılan yaklaşım ölçekli integer kullanımıdır.
    \item PING/PONG, \code{DEOS\_CLASS\_NETWORK / DEOS\_SERVICE\_SYSTEM} altında tanımlıdır.
    \item PING ve PONG command payload'u 4-byte Ping ID içerir ve Length alanı \hex{04}'tür.
    \item PING yalnız unicast olarak geçerlidir; Broadcast PING'e PONG üretilmez.
    \item Bir cevabın Source Node değeri, cevabı üreten mantıksal node kimliği olmalıdır.
\end{itemize}

\section{Açık Konular ve Sonraki ICD Revizyonları}

Mesaj bazında henüz dondurulmamış başlıca CAN-FD protokol alanları şunlardır:

\begin{itemize}
    \item Her mesajın tam command payload byte haritası,
    \item Her komutun minimum ve maksimum fiziksel sınırları,
    \item Her komutun periyodu ve timeout değeri,
    \item Safety mesajlarının nihai priority seviyeleri,
    \item Kalibrasyon payload ayrıntıları,
    \item BMS status mesajının nihai sinyal listesi,
    \item Parameter value type bilgisinin mesaj içinde mi yoksa registry'de mi tutulacağı,
    \item Future firmware update / bootloader mesajları,
    \item Sequence alanının gelecekteki kesin semantiği,
    \item Nominal CAN bitrate ve CAN-FD data-phase bitrate değerleri,
    \item Genel broadcast mesajlarının hangi servis ve komutlarda geçerli olacağı,
    \item PING/PONG timeout ve retry politikaları.
\end{itemize}

Bu alanlar sonraki ICD revizyonlarında ayrıntılandırılacaktır. v1.1'de tanımlanan 29-bit
Identifier yapısı, 4-byte common data header, Length alanı, 60-byte maksimum command payload,
BRS kullanımı ve mantıksal Node ID kuralları sonraki revizyonların referans temelidir.

\end{document}
